% Scribe poster theme (Apache-2.0 / MIT, F. Bruzzone), XeLaTeX
\documentclass[final]{beamer}
\usepackage[orientation=portrait,size=custom,width=84.1,height=118.9,scale=1.13]{beamerposter}
\usetheme{scribeposter}
\usepackage{graphicx,booktabs,tikz,pgfplots}
\pgfplotsset{compat=1.14}
\definecolor{csteal}{RGB}{42,157,176}
\definecolor{csbord}{RGB}{96,0,48}
\definecolor{csgrey}{RGB}{150,150,150}
\newcommand{\placeholder}[2]{\begin{center}\fbox{\begin{minipage}[c][#1][c]{0.92\linewidth}\centering\small #2\end{minipage}}\end{center}}

\title{Can ten qubits finish your sentence?}
\author{Ethan Bandasack \and Malo Leroy \and Thierry Zhang}
\institute{Supervisor: Zeno Toffano \textperiodcentered{} CentraleSupélec, Université Paris-Saclay, L2S}
\date{AI and OneHealth workshop, 20 November 2026}
\posterlogoleft{\includegraphics[height=4cm]{CS}}
\posterlogoright{\includegraphics[height=3.2cm]{L2S_logo}}
\postercontact{Not better than a small classical network}

\begin{document}
\begin{frame}[t]
\vskip0.8cm
\posterbodycenter
\begin{columns}[T, totalwidth=\textwidth]
\begin{column}{0.015\textwidth}\end{column}

\begin{column}{0.30\textwidth}
\begin{block}{Why quantum semantics?}
Quantum-like semantics writes word meanings as vectors and relations between them as operators. Contraries, contradictions and inclusions of the semiotic square become Pauli operators (Galofaro et al., 2018; Dubois and Toffano, 2026). A cosine similarity confuses antonyms, which often share contexts; this vocabulary separates them.
\placeholder{13cm}{Figure: semiotic square with $\sigma_z$ (contraries), $\sigma_x$ (contradictions), $\sigma_y$ (inclusion)}
\end{block}

\begin{block}{Three models, one rule}
Every classical control gets the same inputs, targets and training budget.
\begin{center}\small
\begin{tabular}{lll}
\toprule
Model & Qubits & Readout \\
\midrule
QCSE, reimplemented & 14 & bit-wise $Z$ marginals \\
Hybrid circuit & 8 & classical head \\
Hamiltonian predictor & 10 & Born rule, no head \\
\bottomrule
\end{tabular}
\end{center}
\end{block}
\end{column}

\postercolumnsep

\begin{column}{0.30\textwidth}
\begin{block}{1 \quad A constant guess beats the published metric}
The published score counts a word as right when half of its bits match. A constant prior scores $96.5\%$; QCSE scores $56.1\%$ and recovers no word exactly.
\begin{center}
\begin{tikzpicture}
\begin{axis}[ybar, bar width=14pt, width=\linewidth, height=10cm, ymin=0, ymax=115, ylabel={\%},
  symbolic x coords={half-bit, exact word}, xtick=data, enlarge x limits=0.35,
  legend style={at={(0.5,-0.22)},anchor=north,legend columns=2,draw=none,font=\small},
  nodes near coords, nodes near coords style={font=\small}]
\addplot[fill=csgrey] coordinates {(half-bit,96.5) (exact word,22.8)};
\addplot[fill=csteal] coordinates {(half-bit,56.1) (exact word,0.0)};
\legend{constant prior, QCSE (50 epochs)}
\end{axis}
\end{tikzpicture}
\end{center}
\end{block}

\begin{block}{2 \quad Remove the circuit, and nothing changes}
Bidirectional top-1 accuracy: $50.4\%$ with a trained circuit, $50.0\%$ with the identity circuit, $54.9\%$ with an MLP.
\begin{center}
\begin{tikzpicture}
\begin{axis}[ybar, bar width=18pt, width=\linewidth, height=10cm, ymin=40, ymax=60, ytick={40,50,60}, ylabel={top-1 (\%)},
  symbolic x coords={ZZ circuit, identity, MLP-32}, xtick=data, enlarge x limits=0.25,
  nodes near coords, nodes near coords style={font=\small}]
\addplot[fill=csteal] coordinates {(ZZ circuit,50.38)};
\addplot[fill=csgrey] coordinates {(identity,50.03)};
\addplot[fill=csbord] coordinates {(MLP-32,54.89)};
\end{axis}
\end{tikzpicture}
\end{center}
\end{block}

\begin{block}{3 \quad It learns, but the MLP stays ahead}
On WikiText-2 subwords the quantum NLL falls from $4.28$ to $2.47$ nats; the MLP reaches $2.22$. Coherence helps against a dephased copy, not against a classical network.
\begin{center}
\begin{tikzpicture}
\begin{axis}[width=\linewidth, height=10cm, ylabel={validation NLL (nats)}, xlabel={retained design step},
  xmin=1, xmax=11, xtick={1,6,11}, ymin=2, ymax=4.4, legend style={at={(0.98,0.98)},anchor=north east,font=\small}]
\addplot[thick, csteal, mark=*] coordinates {
  (1,4.2773) (2,3.1975) (3,2.9673) (4,2.9281) (5,2.7870) (6,2.7007)
  (7,2.66776) (8,2.60325) (9,2.57588) (10,2.53554) (11,2.46546)};
\addlegendentry{quantum (Hamiltonian)}
\addplot[thick, csbord, mark=square*] coordinates {
  (1,3.4424) (2,3.2306) (3,2.7966) (4,2.8303) (5,2.4941) (6,2.4393)
  (7,2.41207) (8,2.41207) (9,2.41207) (10,2.41207) (11,2.22277)};
\addlegendentry{classical MLP}
\end{axis}
\end{tikzpicture}
\end{center}
\end{block}
\end{column}

\postercolumnsep

\begin{column}{0.30\textwidth}
\begin{block}{What it means}
\begin{itemize}
\item Most gains came from classical ingredients: learned embeddings, count priors, context pooling.
\item Limits: noiseless simulation, two or three seeds, WikiText test split left unscored.
\item Next: test the semiotic square directly on a relation-classification task.
\end{itemize}
\end{block}

\begin{alertblock}{Take-home}
\begin{itemize}
\item No quantum advantage: every quantum or hybrid model was matched or beaten by a classical model with the same information.
\item Check the metric first: a constant guess scores $96.5\%$ on the published criterion.
\item Always run the identity control.
\end{itemize}
\end{alertblock}

\begin{block}{References}
\small
\begin{itemize}
\item Bengio, Ducharme, Vincent, Jauvin (2003). A neural probabilistic language model. JMLR.
\item Coecke, Sadrzadeh, Clark (2010). Mathematical foundations for a compositional distributional model of meaning.
\item Dubois, Toffano (2026). On quantum relations. Entropy.
\item Galofaro, Toffano, Doan (2018). A quantum-based semiotic model for textual semantics. Kybernetes.
\item Varmantchaonala et al. (2026). QCSE. arXiv:2509.05729.
\end{itemize}
\end{block}
\end{column}

\begin{column}{0.015\textwidth}\end{column}
\end{columns}
\end{frame}
\end{document}
